In the remaining part of this paper we show several examples of polynomial games and recovered their equilibria by the method explained in the \cref{sec:poly}. We implemented the hierarchy of programs \eqref{momsdp} in Julia using \href{https://jump.dev/}{JuMP} \cite{JuMP} and the \href{https://jump.dev/SumOfSquares.jl/stable/}{SumOfSquares.jl} package \cite{PolyJuMP}. We solved the semidefinite programs using \href{https://www.mosek.com/}{Mosek} on a laptop with a 3.1GHz dual-core CPU. Our code is available \href{https://github.com/votroto/PolyNets.jl}{online}.

\begin{example}[\cref{ex.parrilo06.3.1}]
    Consider a two-player zero-sum polynomial game played on the $[-1,1]$ interval. The utility functions are
    \[
        u_{1}(x_{1},x_{2}) = 2x_{1}x_{2}^2 - x_{1}^2 - x_{2} = -u_{2}(x_{1},x_{2}).
    \]
    Since polynomial network games are a generalization of polynomial games, we can use the same framework to solve them. A pure NE is achieved when Player 1 plays $x_{1} \approx 0.4$ and Player~2 plays $x_{2} \approx 0.63$. The optimal payoffs to each player are $\bw \approx (-0.47, 0.47)$. This problem was solved at order 4 of the hierarchy. Due to the low number of variables the problem required only 44 constraints and it was solved within 10\,ms.

    \textcite{parrilo:polynomial} demonstrated that the hierarchy is not necessary for two-player zero-sum games with interval strategy sets. In this case, Lukács' theorem \cite[1.21.1]{szego75orthogonal} states that there exists a Putinar decomposition in terms of rank-1 multipliers such that the term degrees do not exceed that of the input polynomial. The $z\in[-1,1]$ strategy set would be encoded as $1+z \geq 0 \land 1-z \geq 0$ for inputs with odd degrees and $1-z^2 \geq 0$ for even degrees. Both encodings satisfy the Archimedean property as the product of the defining polynomials is positive on a compact set. However, semidefinite solvers usually cannot guarantee rank-1 solutions.
\end{example}

\begin{example}[\cref{ex.kroupa21.5}]\label{ex:separable.network}
    \iffalse
     There are 3 players. All pairs of players are involved in bilateral general-sum games and each player uses the same strategy across all such games. The pairwise utility functions are:
  \begin{equation*}
    \begin{aligned}
      u_{1, 2}(x_1,x_2) &= -2x_1x_2^2 + 5x_1x_2 - x_2\\
      u_{1, 3}(x_1,x_3) &= -2x_1^2 - 4x_1x_3 - 2x_3\\
      u_{2, 1}(x_2,x_1) &= 2x_1x_2^2 - 2x_1^2 - 5x_1x_2 + x_2\\
      u_{2, 3}(x_2,x_3) &= -2x_2x_3^2 - 2x_2^2 + 5x_2x_3\\
      u_{3, 1}(x_3,x_1) &= 4x_1^2 + 4x_1x_3 + 2x_3\\
      u_{3, 2}(x_3,x_2) &= 2x_2x_3^2 + 2x_2^2 - 5x_2x_3
    \end{aligned}
  \end{equation*}
    \fi
    We consider a three-player polynomial network game with a complete graph, where each strategy set is the interval $[-1,1]$ and the utility functions are
    \begin{equation*}\begin{aligned}
        u_{1}(x_{1},x_{2},x_{3}) &= -2x_{1}x_{2}^2 - 2x_{1}^2 + 5x_{1}x_{2} - 4x_{1}x_{3} - x_{2} - 2x_{3},\\
        u_{2}(x_{1},x_{2},x_{3}) &= 2x_{1}x_{2}^2 - 2x_{2}x_{3}^2 - 2x_{1}^2 - 5x_{1}x_{2} - 2x_{2}^2 + 5x_{2}x_{3} + x_{2},\\
        u_{3}(x_{1},x_{2},x_{3}) &= 2x_{2}x_{3}^2 + 4x_{1}^2 + 4x_{1}x_{3} + 2x_{2}^2 - 5x_{2}x_{3} + 2x_{3}.\\
    \end{aligned}\end{equation*}
    The game has the following partially-mixed Nash equilibrium:
    \[
    \begin{aligned}
        \mu^\star_1 &\approx \delta_{-0.06},\\
        \mu^\star_2 &\approx \delta_{0.35},\\
        \mu^\star_3 &\approx 0.72\delta_{1} + 0.28\delta_{-1}.
    \end{aligned}
    \]
    The corresponding utilities are $\bw \approx (-1.22, 0.26, 0.97)$. This problem converged in around $10\,\text{ms}$ at order 4, resulting in 96 constraints.
\end{example}

\begin{example}[\cref{ex.kroupa21.6}]\label{ex:separable.security}
    We consider a polynomial security game with attackers $N_A=\{1,2\}$, defenders $N_D=\{3,4\}$. Player 3 protects targets $t_1$ and $t_2$, while player 4 defends $t_3,t_4$. The capacities of attackers are $a_1=0.5$, $a_2=0.8$, and the defenders' capacities are $d_3=d_4=1$. The efficiency functions are:
    \begin{equation*}\begin{aligned}
        e_{1}(x,y) &= 0.7x^2(1-y^2)\\
        e_{2}(x,y) &= (1.4x - 0.7x^2)(y^2 -2y +1)\\
        e_{3}(x,y) &= (1.8x - 0.9x^2)(1-y^2)\\
        e_{4}(x,y) &= 0.9x^2(y^2 -2y +1)\\
    \end{aligned}\end{equation*}

    Attackers in security games can attack any target, but each defender can only protect targets within their jurisdiction. The interactions can be modelled as a bipartite graph of pairwise general-sum games over targets to which the players allocate resources up to their total capacity. Each target is associated with a resiliency function $e_t$. The total utility of a defender $j$ is \[u_j(x)=\sum_{i\in N_A}\sum_{t\in T_j} e_t(x_i, x_j),\] that is, the sum of resiliency evaluations over all attackers $N_A$ and all targets within their jurisdiction $T_j$; the total utility of an attacker $i$ is \[u_i(x)=a_i - \sum_{j\in N_D}\sum_{t\in T_j} e_t(x_i, x_j).\] Thus the games are globally constant sum.

    Our algorithm returned the following partially-mixed strategy profile at order 6:
    \[
    \begin{aligned}
    \mu^\star_1 &\approx \delta_{(0.29, 0, 0, 0.21)}, \\
    \mu^\star_2 &\approx \delta_{(0.36, 0.13, 0.04, 0.24)}, \\
    \mu^\star_3 &\approx 0.46\delta_{(0.33, 0.67)} + 0.54\delta_{(0.69, 0.31)}, \\
    \mu^\star_4 &\approx 0.02\delta_{(0.46, 0.54)} + 0.06\delta_{(0.08, 0.92)} + 0.92\delta_{(0.91, 0.09)}
    \end{aligned}
    \]
    The payoffs to each player are $\bw \approx (0.1, 0.3, -0.17, -0.23)$. While the payoffs converged at order 4 already and took only 3 seconds to solve, the corresponding strategies could not be extracted until order 6, which took 6 minutes to solve. Due to the high dimensionality of the utility functions and the fast growth of the hierarchy, the resulting SDP contained over 25 thousand constraints.
\end{example}
